% !TeX program = xelatex
\documentclass[journal,wp,nowatermark,nobib,clean]{paper}
\usepackage{longtable,needspace,listings,tikz,xurl,placeins}
\usetikzlibrary{arrows.meta,positioning,fit,backgrounds,calc}
% Preserve the kit class; all project-specific styling stays in this file.
\definecolor{ink}{HTML}{211F1B}
\definecolor{pine}{HTML}{315F5A}
\definecolor{gold}{HTML}{B27D45}
\definecolor{mist}{HTML}{F1F3EF}
\definecolor{rulegray}{HTML}{BBC2BB}
\hypersetup{pdftitle={TROPOS 3.2.9 自适应银行智能体技术文档},pdfauthor={咪咪大王队：刘水、余冠辰},pdfsubject={2026 FinTechathon 国内人工智能赛道 技术设计与实现},bookmarksnumbered=true,hypertexnames=false}
\doctype{技术文档}
\assigntitle{TROPOS 自适应银行智能体}
\assignsubtitle{系统架构与关键技术说明}
\authorname{\shortstack[l]{咪咪大王队\\刘水、余冠辰（复旦大学经济学院）}}
\affiliation{TROPOS-TD-3.2.9}
\submitdate{2026 年 11 月 2 日}
\wpshorttitle{TROPOS 3.2.9 技术设计与实现}
\fancyhead[C]{\footnotesize\paperkaiti\color{wprust}TROPOS 3.2.9\quad 技术设计与实现}
\renewcommand{\arraystretch}{1.22}
\renewcommand{\paperlogofile}{assets/brand-logotype.pdf}
\renewcommand{\papernamelabel}{团队}
\renewcommand{\paperidlabel}{编号}
\renewcommand{\papercourselabel}{赛题}
\renewcommand{\paperdatelabel}{日期}
\course{AI 代人操作银行（安全与体验）}
\wpnumber{技术文档\quad 系统版本 3.2.9\quad 修订版 2.0}
\newcommand{\contestheading}{%
  {\small\paperkaiti\color{wprust}
  2026（第八届）FinTechathon 深圳国际金融科技大赛\\[.3em]
  微众银行金融科技大学生挑战赛\quad 国内人工智能赛道\par}\vspace{.8cm}}
\patchcmd{\makecover}{\vspace*{0.6cm}}{\contestheading\vspace*{0.6cm}}{}{\PackageError{tropos}{Cover patch failed}{}}

\setcounter{tocdepth}{2}
\setcounter{secnumdepth}{2}
\emergencystretch=2em
\raggedbottom
\widowpenalty=10000\clubpenalty=10000\displaywidowpenalty=10000\postdisplaypenalty=10000
\AddToHook{cmd/subsection/before}{\FloatBarrier}
\setlength{\headheight}{15pt}
\setlist[itemize]{itemsep=.18em,topsep=.35em,parsep=0pt}
\setlist[enumerate]{itemsep=.2em,topsep=.35em,parsep=0pt}
\newcommand{\code}[1]{{\small\nolinkurl{#1}}}
\newcommand{\source}[1]{\par\nopagebreak\vspace{.25em}{\footnotesize\color{wprust}实现依据：\code{#1}。\par}}
\renewcommand{\tabnote}[1]{\par\vspace{.25em}{\footnotesize #1\par}}
\newcommand{\sectionstart}{\FloatBarrier\par\Needspace{14\baselineskip}}
\newcolumntype{L}[1]{>{\raggedright\arraybackslash}p{#1}}
\newcolumntype{Y}{>{\raggedright\arraybackslash}X}
\lstset{basicstyle=\ttfamily\footnotesize,breaklines=true,columns=fullflexible,keepspaces=true,showstringspaces=false,frame=none,aboveskip=.6em,belowskip=.6em,xleftmargin=.7em,backgroundcolor=\color{white}}
\tikzset{box/.style={draw=wpblue,fill=wpcream,rounded corners=1pt,align=center,font=\small,inner sep=7pt,text=ink},arr/.style={-{Latex[length=2mm]},draw=wprust,thick},domain/.style={draw=wprust,dashed,rounded corners=3pt,inner sep=9pt}}
\begin{document}
\special{pdf:minorversion 7}
\disablewatermark
\makecover
\setcounter{page}{1}
\makewptitle
\begin{abstract}
\hspace*{2em}TROPOS 面向自然语言办理银行事务，研究如何让智能体理解人的财务处境，并在可核验的约束下形成安排。项目将长期偏好与短期现金状态分别记录，通过知己、善谋、择时、守诺四种职责，连接画像、策略、计划、执行与反馈。相同需求因现金缓冲、到期义务与目标期限而产生不同安排；收入或账单变化后，未执行计划需要重新核验。

当前 3.2.9 实现合成银行沙箱中的财务体检、转账、订阅办理、卡片及理财演示、多步骤任务与生活规划。模型参与语义理解、候选规划和只读回答；金额、授权与记账事实由确定性服务端控制。银行账本与银行回执为合成数据；美团官方连接器可能涉及真实授权、地址和订单，实际支付在美团完成。系统未接入真实银行，也没有后台自动资金执行器。

本文说明系统架构、现金算法、确认状态机、数据与外部服务边界，并提供本次复跑结果及源码指纹，供评审与复现。
\end{abstract}
\vspace{.8em}
\noindent{\small\textbf{关键词}\quad 动态金融画像；现金流约束；分级授权；可解释规划；幂等记账。\par}
\vspace{1.2em}
\noindent{\small\textbf{版本基线}\quad 后端 \code{2a33063}；界面采用 2026 年 11 月 2 日本地工作区水墨风更新。验证结果归属于测试运行时快照，详见第~\ref{sec:validation}~节与随附证据。\par}
\clearpage
\maketoc
\section{项目出发点与设计目标}
\label{sec:positioning}
\subsection{从一笔指令到一份适合当下的安排}
银行提供的操作能力日益丰富，但用户仍要在收入、账单、家庭责任与生活目标之间，自己拼出一份安排。“帮我打理这两万元”对现金紧张者和储备充足者，不应产生相同的行动；昨天能够执行的计划，也可能因工资延期而在今天失去可行性。TROPOS 的出发点，是让银行智能体先理解资金在生活中的位置，再决定可以准备什么操作。

\begin{keybox}[核心设计]
以动态画像形成安排，以确定性内核守住执行边界。模型提出候选，用户保留决定权，账本事实由服务端核验。
\end{keybox}

“千人千策，随境而变”在工程上对应两个要求。一是个体差异进入决策条件：现金底线、必要支付、期限与用户主动声明的偏好，影响建议空间和核验结果。二是状态变化进入后续决策：收入到账、延期、账单变化或画像更正，会改变账户状态或版本，使旧方案接受重新检查。系统不把一次问卷变成永久标签，也不因为一次收入冲击擅自改变用户的长期目标。

项目面向比赛评审和可复现的技术验证，核心问题是如何把模型生成与金融状态变更衔接起来：允许模型理解自然语言与提出候选安排，同时以确定性业务服务掌握事实与执行权。本文将产品立意落实到数据结构、约束、确认过程和回执，而不是以功能数量代替这一主线。

\subsection{四种资金需要与项目版图}
Protect、Recover、Accumulate、Grow 分别对应现金流守护、负债减压、财富积累与长期资金安排。它们是随处境变化的任务重点，可以并存，也可以调整；不能依据年龄、职业等身份直接为用户分配固定金融策略。

\begin{table}[htbp]\centering\small
\caption{资金需要与当前实现的对应关系}
\begin{tabularx}{\textwidth}{@{}L{2.2cm}YY@{}}\toprule
任务重点 & 要回答的问题 & 当前可体验的支撑 \\\midrule
现金流守护 & 必要账单和生活缓冲能否保留 & 30 天现金路径、现金底线、事件冲击复核 \\
负债减压 & 到期义务怎样排序与核对 & 还款与账单汇总、订阅复核、多步骤任务 \\
财富积累 & 日常之外能为目标留出多少空间 & 目标储蓄草案、生活安排比较、预算核验 \\
长期资金安排 & 用钱期限与风险约束是否匹配 & 主动风险声明、期限与流动性核对、理财沙箱 \\\bottomrule
\end{tabularx}
\end{table}

产品由三层场景连接起来：银行事务提供转账、查询、订阅及卡片等操作工具；现金规划判断近期约束与目标空间；生活联动把家庭、购物、教育、健康和出行等需要放在同一份资金日历中。长期愿景是持续适应人的财务变化；当前落点仍是有限期限、明确授权、可核验的比赛原型。

\subsection{赛题对应与实现边界}
赛事国内人工智能方向聚焦银行 AI 智能体，通过自然语言理解意图、规划并办理业务，兼顾安全与体验。TROPOS 以转账、账单分析与订阅办理为基础，扩展卡片、理财沙箱及生活联动。赛题名称依据赛事官网，访问日期为 2026 年 11 月 2 日\footnote{FinTechathon 赛事官网：\url{https://fintechathon.g-ican.com/}。本文不据官网介绍推定未公开的评分权重或宣称已经通过赛事验收。}。

银行沙箱中的账户、余额、交易、复式分录和银行回执均为合成记录，不连接真实银行清算。美团官方链路、外部模型和 SMTP 则可能产生真实的对外请求；它们分别受连接器、数据边界与部署配置约束。真实美团订单创建不等于合成银行付款，跳转收银台也不等于支付成功。

系统当前没有实名核验、银行卡要素认证、生产清算、持续后台自动付款或经过验证的完整生命周期模型。页面中的资金预测是基于已知输入的条件计算，不是工资到账保证、收益承诺或完整财务风险评估。

\sectionstart
\section{总体架构与四引擎协作}
\label{sec:architecture}
\subsection{分层结构与信任边界}
系统采用同源网页、Python 标准库 HTTP 服务与 SQLite 数据库。页面负责输入、展示、选择和独立确认；网关解析会话及请求；业务服务读取事实、计算约束并执行状态变更。模型返回值与外部网页内容均作为待检查的数据，不因来源于模型而获得授权。

\begin{figure}[htbp]\centering
\caption{系统分层与对外连接边界}
\begin{tikzpicture}[node distance=7mm and 6mm]
\node[box,text width=12cm] (ui) {同源用户界面\quad 首页、工作台、财务体检、生活规划、模型中心};
\node[box,text width=12cm,below=of ui] (http) {HTTP 接入\quad 路由与请求校验、用户会话、读取同意};
\node[box,text width=5.5cm,below left=9mm and -5.65cm of http] (model) {受限智能编排\\意图、澄清、候选方案、只读回答};
\node[box,text width=5.5cm,right=6mm of model] (core) {确定性业务内核\\现金计算、版本复核、授权与记账};
\node[box,text width=5.5cm,below=of model] (ext) {外部服务\\模型与搜索、SMTP、美团连接器};
\node[box,text width=5.5cm,below=of core] (db) {SQLite 与独立运行状态\\合成账本、回执、审计、用户分区};
\draw[arr] (ui)--(http);\draw[arr] (http.south)-|(model.north);\draw[arr] (http.south)-|(core.north);
\draw[arr] (model)--node[above,font=\scriptsize]{校验候选}(core);
\draw[arr] (model)--(ext);\draw[arr] (core)--(db);
\end{tikzpicture}
\tabnote{外部连接不共享银行账本执行权。美团真实支付在用户与美团之间完成，订单状态与合成银行交易分别记录。}
\end{figure}

后端以模块而非独立模型进程组织职责：\code{server.py} 提供接入，\code{service.py} 管理账户、身份和账本，\code{conversation.py} 与 \code{agent_ops.py} 编排任务，\code{model.py} 适配模型，\code{life_planner.py} 执行有限期规划。四引擎是贯穿这些模块的逻辑职责，不要求四个大模型分别驻留，也不暗示存在未实现的自主协商网络。

\subsection{职责到工程模块的映射}
\begin{figure}[htbp]\centering
\caption{四种职责围绕同一份用户状态协作}
\begin{tikzpicture}[every node/.style={text=wpblue}]
\foreach \x/\g/\en/\cn in {0/知/ETHOS/理解处境,4.1/谋/METIS/形成安排,8.2/时/KAIROS/复核变化}{
  \node[draw=wpblue,circle,fill=wpcream,minimum size=15mm,line width=.8pt] (r\en) at (\x,0) {\Large\paperkaiti \g};
  \node[align=center,font=\small] at (\x,-1.22) {\en\\\cn};
}
\draw[-{Latex[length=2mm]},wpblue] (rETHOS.east)--(rMETIS.west);
\draw[-{Latex[length=2mm]},wpblue] (rMETIS.east)--(rKAIROS.west);
\node[draw=wprust,circle,fill=wpcream,minimum size=15mm,line width=.8pt,text=wprust] (guard) at (4.1,-3.1) {\Large\paperkaiti 守};
\node[align=center,font=\small,text=wprust] at (4.1,-4.3) {ARGOS 守诺\\授权、金额、版本与回执的跨层核验};
\draw[wprust,dashed] (0,-1.65)--(0,-1.85)--(8.2,-1.85)--(8.2,-1.65);
\draw[wprust,dashed] (guard.north)--(4.1,-1.85);
\end{tikzpicture}
\fignote{四者表示逻辑职责，并非四个独立大模型。守诺贯穿理解、规划和执行，不排在最后才介入。}
\end{figure}

\paragraph{ETHOS 知己}
以用户目标、期限、风险声明、当前现金与更正信息形成版本化画像，并解释当前任务侧重。实现位于画像规则及 \code{service.py}，长期声明与账本状态分别保存。

\paragraph{METIS 善谋}
结合账单、现金约束和允许的行动，生成候选方案与有序任务。实现分布在 \code{model.py}、\code{life_planner.py} 与 \code{agent_ops.py}；候选必须经过确定性核验。

\paragraph{KAIROS 择时}
检查收入事件、日期条件、有效期与版本变化，产出等待、失效和重新生成的状态。实现位于 \code{service.py}、\code{life_service.py} 与对话恢复流程。

\paragraph{ARGOS 守诺}
核对身份、授权、对象、金额、版本和回执，必要时拒绝推进。接入、业务服务和任务执行共同承担这项职责，最终状态变更受服务端约束。


ARGOS 的作用贯穿生成与执行：草案阶段检查可行性，确认阶段复核当前事实，执行阶段检查幂等与账本一致性。它拥有停止推进的业务效果，而不是在流程末尾补一句安全提示。KAIROS 的现有触发来自明确请求、演示事件与确认时检查；没有后台常驻事件执行器，条件满足后仍须重新核验与确认。

\subsection{画像到反馈的闭环}
用户首先声明目标、现金底线、用钱期限、流动性与风险意愿，并明确是否同意读取所需信息。画像层保留这些长期声明，账本层单独记录余额、账单、工资状态和版本。策略层据此约束候选空间，计划层形成可以逐项查看的步骤，业务层只执行当前有效且已经确认的动作。回执、用户更正和新事件成为后续判断的输入。

例如工资延期只改变短期收入时点，不等于用户同意提高风险或降低现金底线。现金路径重算后，原先可行的转账草案可能因账户版本变化而失效。系统应解释为何需要重看、保留已完成操作的回执，并为未执行部分重新准备草案；不把原来的授权延伸为新安排的空白授权。

\sectionstart
\section{确定性现金计算与可解释规划}
\subsection{现金路径的精确定义}
金融计算使用整数分。金额入口拒绝布尔值、非有限数和超过分精度的输入，避免把浮点累计误差带入余额判断。设主账户已到账余额为 $B$，假设立即转出额为 $D$，每日演示开销为 $d$，第 $t$ 天的有符号事件金额之和为 $F_t$。则：
\begin{align}
C_0&=B-D,\\
C_t&=C_{t-1}-d+F_t,\qquad t=1,\ldots,30,\\
C_{\min}&=\min_{0\leq t\leq30}C_t,\qquad
G=\max\{0,R-C_{\min}\}.
\end{align}
其中 $R$ 为用户现金底线，$G$ 为预测底线缺口。路径包含今日和未来 30 天，共 31 个节点；最低值并列时取首次出现日。\code{free} 定义为 $C_{\min}-R$，可以为负值，不应直接称为可花余额。

事件包括既有账单、已登记的额外义务、有效订阅，以及尚未到账的预计工资。收到工资事件后，当前余额增加，同时移除预测中的这笔工资，防止重复计薪。延期改变预计到账日，移出观察窗口的收入不再计入 30 天路径。尚未确认的转账计划不是实际义务，不能因为模型提出建议就写入现金事实。

代码按每日开销和当日事件净额计算，只保存日终节点，不模拟同日先扣款后收入造成的日内缺口。第 0 天事件不在未来日循环中执行。当前演示账本与财务体检采用固定基准日，不能把输出称为已接入真实流水的滚动月报。
\source{backend/service.py: BankService._cash; backend/insights.py: build_insights}

\subsection{一组可以独立复算的合成算例}
设今日余额为 10,000 元，现金底线为 4,000 元，每日支出 100 元，第 5 天支付账单 3,000 元，第 10 天预计工资 6,000 元，第 12 天订阅扣费 100 元。以下均为说明算法而构造的输入，不是用户真实账本。

\begin{table}[htbp]\centering\small
\caption{同一账本在不同条件下的现金最低点}
\begin{tabularx}{\textwidth}{@{}YrrY@{}}\toprule
条件 & 最低现金（元） & 首次最低日 & 与底线比较 \\\midrule
原条件，不转出 & 6,100 & 第 9 天 & 高于底线 2,100 元 \\
原条件，立即转出 2,000 元 & 4,100 & 第 9 天 & 高于底线 100 元 \\
原条件，立即转出 2,500 元 & 3,600 & 第 9 天 & 低于底线 400 元 \\
工资延期 15 天，不转出 & 4,500 & 第 24 天 & 高于底线 500 元 \\
延期后，仍转出 2,000 元 & 2,500 & 第 24 天 & 低于底线 1,500 元 \\\bottomrule
\end{tabularx}
\end{table}
原条件下第 9 天余额为 $10000-9\times100-3000=6100$ 元；延期后第 24 天为 $10000-24\times100-3000-100=4500$ 元。相同转账额在两种条件下得到不同核验结果，体现状态变化对行动的约束。通过现金检查只说明这一组输入下的可行性，不产生用户授权，也不保证未来输入准确。

\subsection{财务体检的证据与规则}
体检复用服务端快照和现金路径，展示现金压力、小额支出、订阅复核项及其依据。单笔不超过 100 元且不属于内部调拨或还款的可见出账用于小额支出汇总；达到 2,000 元的出账或包含指令型文本的备注提示本人复核。低使用和重复订阅标记来自演示资料，不由模型假装推断真实使用频率。

现金低于用户底线时标记 \code{urgent}；没有现金缺口但存在复核信号时为 \code{attention}；没有上述信号时为 \code{steady}。这些是可解释规则状态，不是信用评分、欺诈判断或经过校准的风险概率。

取消单项订阅的比较，从窗口内对应续费日开始加回费用，不增加今日余额。由于最低点可能移动，多项取消对“最低余额改善”的影响不能把单项结果简单相加。体检同时列明数据、规则、未知项和决策边界；金融白话解释是只读知识入口，长辈模式只调整显示与关键确认信息的可读性。

\subsection{生活规划的有限离散搜索}
生活场景由用户声明金额、首选日、最晚日、重要性，以及是否允许替代或跳过。系统对保留原安排、缓冲优先、兼顾生活三种方案，比较基准、收入额外延期、突发支出、二者叠加四类压力条件。各情景没有估计概率；工资已到账时，延期情景可能与基准重合。

兼顾方案使用宽度为 128 的束搜索。候选日期来自首选日、最晚日以及允许窗口内的预计工资后日期，不逐日穷举；金额只保留用户原值、主动允许的替代金额，或在允许且非必要时跳过。必要场景不能自动取消。排序依次比较四情景下最坏支付缺口、最坏现金底线缺口、生活调整成本、节点缺口累计和日期指标；最终保留原安排与缓冲方案作为候选。

调整成本采用必要 100、重要 30、可调整 5 的演示权重，服务于可解释排序，不是估计效用或货币福利。有限束宽不保证全局最优。不动用资产只展示，不计入当前支付能力；缓冲敏感性固定同一方案与目标额，只改变初始现金假设。

采用生活安排会在事务中重算、复核双版本并登记未来义务，当前余额不扣款；采用不意味着所有压力情景均无风险。目标储蓄需另行生成转账草案，确认时走银行网关。生活分析、登记义务与实际付款具有不同状态，界面与回执必须分别表达。
\source{backend/life_planner.py: analyze, _choices, _rank, _policy; backend/life_service.py: commit_plan, propose_goal}

\sectionstart
\section{对话办理与状态变更控制}
\subsection{从自然语言到待确认对象}
对话层优先处理明确安全门与当前业务上下文，再选择本地规则、结构化理解、财务方案或普通只读回答。歧义收款人、缺失金额及不明确对象进入澄清，不能由模型凭空补齐。金额更正会创建新转账草案并撤销旧草案，已结算记录不能被聊天改写成未执行状态。

聊天中的“确认”展示确认入口，不直接执行转账。实际确认使用业务对象、精确令牌和明确布尔确认。对话并非完全只读：创建草案、写会话、撤销待办等可以改变服务端状态，但它们与资金记账不同。

\begin{figure}[htbp]\centering
\caption{普通转账从草案到回执的核验链}
\begin{tikzpicture}[node distance=6mm]
\node[box,text width=11.8cm] (a) {当前请求与澄清\quad 不把历史文字当作资金授权};
\node[box,text width=11.8cm,below=of a] (b) {草案生成\quad 登记收款人、金额、余额、现金路径与前置条件};
\node[box,text width=11.8cm,below=of b] (c) {用户独立确认\quad 对象、金额、令牌与有效窗口};
\node[box,text width=11.8cm,below=of c] (d) {写事务内复核\quad 用户归属、版本、过期、最新条件、原交易};
\node[box,text width=11.8cm,below=of d] (e) {首次有效请求：原子记账\qquad 重试：返回原回执};
\foreach \x/\y in {a/b,b/c,c/d,d/e}\draw[arr](\x)--(\y);
\end{tikzpicture}
\tabnote{不通过的对象保持等待、阻断、失效或撤销状态；不能沿图示直通执行。}
\end{figure}

\subsection{普通转账的草案与确认契约}
普通转账草案核验登记收款人、合法金额、单笔不超过 50,000 元、当前已到账余额、转出后现金底线、未来现金路径和工资前置条件。检查通过时为 \code{pending_confirmation}；工资条件未满足时可以显示 \code{waiting_condition}，但该状态不表示其他检查已经通过。其余不可行条件为 \code{blocked}。

草案绑定用户、账户修订号、随机确认令牌和 600 秒有效期。确认要求 \code{confirmation} 严格为布尔值 true、读取同意、精确令牌以及当前账户版本；与生活目标关联的草案还检查生活版本与无待比较场景等条件。最终普通现金检查与记账位于同一写事务中。

\begin{table}[htbp]\centering\small
\caption{普通转账状态的存储与展示口径}
\begin{tabularx}{\textwidth}{@{}L{4.0cm}Y@{}}\toprule
状态 & 含义及后续处理 \\\midrule
\code{pending_confirmation} & 草案生成时可行；确认时仍须复核 \\
\code{waiting_condition} & 工资等条件未满足；满足后重新准备或复核，不自动付款 \\
\code{blocked} & 当前金额、余额、现金路径或条件不满足 \\
\code{expired} / \code{stale} & 窗口过期，或账户、生活计划发生变化；重建有效草案 \\
\code{revoked} & 用户已撤销尚未结算的草案 \\
\code{settled} / \code{executed} & 数据存储记为 settled；已有交易时展示为 executed，查询原回执 \\\bottomrule
\end{tabularx}
\end{table}

部分状态是读取时投影，不是均以同名字段写回数据库。已结算交易优先投影为已执行；不得仅从展示字符串推断底层状态迁移。已记账交易不能通过撤销草案恢复到未执行。

事务以方案 ID 作为规范交易幂等标识；相同用户和方案已有交易时，在核验原确认信息后返回原回执。更换客户端请求键不会因此产生第二笔扣款。首次记账更新余额、两条方向相反的分录、交易、授权、审计与方案状态，并递增账户修订号。本地事务与唯一约束保证的是沙箱账本内的一致性，不能推广为外部支付系统的“恰好一次”承诺。
\source{backend/service.py: _propose_transfer, _present_proposal, confirm, _post, revoke}

\subsection{多步骤任务与部分完成}
复合请求最多处理八个分句，步骤分别标记只读完成、可确认、阻断、需澄清或拒绝。首次整单确认要求任务令牌、账户修订号、600 秒有效期和明确批准的步骤集合，只能批准当前可执行的步骤；挂失另需卡号后四位，但这不等于生产强身份认证。

执行按顺序推进：未批准的步骤记为跳过，成功步骤保留回执，业务失败步骤记为失败，之后获批准的步骤暂停。正常返回的任务终态 \code{executed} 表示整单执行尝试已经结束，可能包含全部跳过或部分失败，不能译为“全部成功”。重复确认返回原结果，跳过和暂停项不能在同一已结束任务中直接补做，应核对原结果后重新规划。

业务失败可以保留此前成功结果；未捕获异常则由外层事务回滚。转账与申购按前步骤更新后的账本核对资金。普通转账的 50,000 元限额不能未经核查推广为所有任务、所有入口或每日统一限额；“剩下的”金额分配还存在独立的额度计算规则。

当前不同变更入口的确认契约并不统一。单项订阅取消依靠显式确认、订阅状态与既有回执处理重试，没有普通转账草案的全部令牌、TTL 和调用方版本字段。卡片及理财任务的执行分支，也未逐一重复规划阶段的所有条件检查。后续应统一“对象、版本、时限、最新条件、幂等”的动作契约；本文不把局部保证写成全局保证。
\source{backend/agent_ops.py: plan, confirm, _execute; backend/service.py: _cancel_subscription}

\sectionstart
\section{模型中心与受限智能编排}
\subsection{模型配置与本地回退}
模型中心区分本地规则、平台模型和用户自带模型。平台密钥来自部署配置或可信本机管理员的运行时配置；个人 API Key 按认证用户隔离，仅存服务进程内存，重启后清除。选择状态属于用户自身，不应改变其他用户的模型选择。服务未配置、配置失效或输出不可用时，系统采用本地规则或明确失败说明。

\begin{table}[htbp]\centering\small
\caption{代码中的供应商适配与使用条件}
\begin{tabularx}{\textwidth}{@{}L{1.8cm}YY@{}}\toprule
服务 & 适配路径 & 应核验的运行条件 \\\midrule
DeepSeek & 对话接口与检索适配 & 模型标识、凭据、当前检索入口与可用权限 \\
OpenAI & Responses 与原生网页检索适配 & 选用模型及其工具支持、引用返回情况 \\
Kimi & 兼容对话与 Search Pro 适配 & 凭据、配额、模型与检索返回格式 \\
千问 & 兼容对话及供应商专用能力适配 & 具体模型、思考开关、检索和输出方式 \\
自定义 & OpenAI 兼容入口 & 配置地址与服务端真实支持范围 \\
本地规则 & 确定性解析与业务规则 & 无外部模型依赖，仍需用户确认写操作 \\\bottomrule
\end{tabularx}
\tabnote{该表描述项目代码适配，不是供应商全量能力清单，也不代表本次逐一完成真实联网验收。}
\end{table}

\subsection{发往模型的数据契约}
普通意图理解主要使用当前用户消息；只读多轮问答可使用最近几轮文字帮助理解指代；生活或财务方案在同意读取画像后，使用服务端构造的财务汇总。汇总可包含余额、现金底线、现金路径、分类和必要账单合计，不主动附加原始流水备注、商户明细、收款账号、用户标识或会话令牌。

这项限制针对服务端构造的数据字段。用户自己输入或历史消息中可能已经含有姓名、账号、地址或其他敏感内容，不能据此承诺“任何敏感信息都不会发送”。API Key 不进入提示正文和响应回显，但用于所选供应商请求的鉴权。用户自带地址的信任与数据接收方应在选择模型时说明，不能把兼容接口等同于可信服务。

\subsection{结构化方案与服务端上限}
模型方案只是候选。服务端检查字段集合、重复 JSON 键、类型、有限数值、文字长度、金额合计、服务入口白名单与来源编号，并拒绝把尚未执行的动作写成已完成事实的输出。来源材料作为不可信证据使用，网页不能赋予付款权限；结果缺少可用来源时不包装成已经核实的联网结论。

财务方案的预算由服务端先计算上限，模型只能在上限内提出安排。当前确定性建议路径的上限为
\begin{equation}
U=\max\bigl\{0,\min(P,\ B-R,\ C_{\min}-R)\bigr\},
\end{equation}
其中收入稳定且无演示冲击时 $P=2000$ 元，否则为 1000 元。这是当前非必要支出建议的演示策略上限，不是适用于所有转账的监管限额。生活规划的目标上限则取请求目标额、四压力情景的最低现金减底线、当前余额减底线三者的非负最小值。校验应使用对应接口返回的上限，不能以一个通用公式替代所有业务规则。

涉及金额的模型文本不能替代确定性计算，关键词或语义规则也不能证明所有幻觉都能被识别。当前实现依靠受限字段、动作白名单和最终执行网关控制实际副作用；自然语言错误仍须以准确的界面状态和回执予以纠正。

\subsection{3.2.9 的千问稳定性修复}
3.2.9 将千问普通自由对话调整为单次非流式请求，固定关闭该路径的 \code{enable_thinking}，减少隐藏思考消耗输出额度却没有可见回答的情况。连接测试采用一次最小只读对话，不再先理解再自然语言生成；上游失败后转本地只读规则，不在同轮悄悄追加兼容重试。

浏览器对只读对话、模型连接测试与交易类请求分别提示超时。只读超时不提示用户查询不存在的“原交易”；真实变更结果不明时则依据原凭证查询。修复只适用于已经调整的路径，不能推定所有模型、所有接口都固定非流式或关闭思考。
\source{backend/model.py; backend/model_center.py; assets/backend-client.js}

\sectionstart
\section{购物链路与外部服务隔离}
\subsection{模拟购物与真实跑腿分开记账}
Mock 链路覆盖合成商品搜索、预算核对、逐项确认、模拟扣款与模拟履约，状态如 \code{mock_paid} 只代表沙箱记录。官方连接器则通过用户授权与美团组件，读取真实地址、询价、准备帮买订单，确认后调用供应商，再由用户进入美团收银台完成真实支付。通用商品搜索仍使用 Mock 数据，不能据跑腿帮买宣称真实餐厅菜单外卖已打通。

\begin{table}[htbp]\centering\small
\caption{两种购物链路的事实边界}
\begin{tabularx}{\textwidth}{@{}L{2.15cm}YY@{}}\toprule
比较项 & Mock 购物 & 美团官方跑腿连接器 \\\midrule
信息来源 & 合成商品与商户 & 授权地址、询价及供应商响应 \\
确认后动作 & 沙箱订单与合成扣款 & 真实帮买订单提交及支付跳转 \\
支付事实 & 仅模拟支付 & 由用户在美团侧完成，须查询真实状态 \\
银行账本 & 写入合成交易与回执 & 不扣合成银行余额 \\
失败处理 & 本地事务与原订单查询 & 保留未知状态、按原凭证核对，不盲目重提 \\\bottomrule
\end{tabularx}
\end{table}

\subsection{用户授权与提交控制}
官方链路须具备 Node 程序、Skill、授权助手和状态目录。每个用户先建立供应商授权，预览绑定用户、地址与询价条件。确认时检查对象归属、令牌、账户版本、过期和预算；较高配送费需要额外确认。订单创建只说明供应商接受了请求，不代表已经支付、接单或履约。

提交超时或结果不明时保留 \code{unknown}，限制同一询价直接重提。调用方应查询原订单状态并解释不确定性，不能因为没有及时收到响应便假定未下单。跨系统幂等能力取决于供应商契约，本地请求去重不能替代对方的执行保证。

真实配送必须向供应商提供收件姓名、电话、地址或坐标，以及物品要求、预算和备注。这是功能所必需的数据传输，与模型汇总字段裁剪是两种不同的数据边界。浏览器展示脱敏地址，不意味着供应商提交报文也只含脱敏信息。

\subsection{组件调用与凭据空间}
连接器通过参数数组启动 Node 子进程，不经 shell 拼接命令。授权缓存目录使用用户 ID 的 SHA-256 派生路径，子进程 HOME、临时目录与状态缓存分开配置；Linux 路径权限设置限制本机普通访问。该设计隔离凭据命名空间，不能等同于独立容器或操作系统级沙箱。

访问 Token 可能由官方授权助手缓存到文件，并在调用时进入子进程环境；部分询价凭证和供应商订单对象会保存到数据库。因此应区分模型密钥、供应商授权 Token、业务订单凭证和公开状态，不使用“所有第三方凭据都不落盘”这样的概括。

地址接口可能出现进程退出码为零但业务失败的情况，系统还需检查供应商业务响应。临时手动地址提供有限兜底，带 HMAC 签名的编码凭证用于完整性验证；签名不等于加密，不能把其中可解码内容当成保密存储。移动端提供打开授权页、手动刷新、重试与重新授权入口。
\source{backend/shopping.py; backend/meituan_paotui.py; scripts/package_meituan_connector.py}

\subsection{邮件服务的真实副作用}
配置 SMTP 时，邮箱验证与密码重置会向真实邮箱发送链接；未配置时，邮件落在权限受限的本机 outbox，属于开发行为，不表示用户已收到邮件。发送过程会向邮件服务提供收件地址。传输是否启用 STARTTLS 取决于配置，不能把可关闭的选项描述为任何环境下都已强制加密。
\source{backend/mailer.py}

\sectionstart
\section{数据模型与身份安全}
\subsection{持久化对象与版本}
SQLite 保存用户、画像、合成账户、交易、分录、幂等信息、方案、事件、对话、生活计划和购物状态。资源查询附带当前会话解析得到的用户 ID，不能由客户端请求字段任意替换归属。写事务采用 \texttt{BEGIN IMMEDIATE}，并用唯一约束保证用户范围内的请求去重。

\begin{table}[htbp]\centering\small
\caption{主要持久化对象及用途}
\begin{tabularx}{\textwidth}{@{}L{2.2cm}L{4.25cm}Y@{}}\toprule
对象组 & 代表表 & 主要约束或版本 \\\midrule
用户与画像 & \code{users}, \code{profiles}, \code{profile_updates} & 账户 revision、画像版本、授权读取状态 \\
账本与回执 & \code{accounts}, \code{transactions}, \code{entries} & 整数分、交易唯一键、相反方向分录 \\
请求与方案 & \code{idempotency}, \code{proposals}, \code{authorizations} & 请求指纹、令牌、有效期与用户归属 \\
身份与挑战 & \code{auth_accounts}, \code{sessions}, \code{auth_challenges} & 口令及令牌摘要、消费状态与过期 \\
多因素认证 & \code{mfa_settings}, \code{mfa_recovery_codes} & TOTP 计数、防重放、恢复码单次消费 \\
对话与任务 & \code{conversations}, \code{conversation_requests}, \code{missions} & 对话版本、资源引用、步骤执行结果 \\
生活与购物 & \code{life_*}, \code{shopping_*} & 生活版本、账本版本、订单与供应商绑定 \\\bottomrule
\end{tabularx}
\end{table}

账户修订号用于发现财务事实变化；生活版本用于发现方案输入变化；对话版本用于并发消息与上下文一致性。三者并不互相替代。对话请求指纹绑定用户、消息与原会话标识，同键异内容应拒绝；重放可以重新读取业务资源状态，而不是把历史文本里的完成描述当作最新事实。

\subsection{注册账户与共享演示账户}
注册账户生成独立用户 ID 和合成账本。密码使用独立盐值的 scrypt 慢哈希；会话令牌使用高熵随机值，数据库保存摘要，有效期为 24 小时。连续五次错误口令触发 15 分钟限制；重置密码撤销已有会话。验证与重置链接使用短期、单次消费挑战，邮箱未验证并不阻止当前注册流程立即签发会话，不能把邮箱验证称为所有访问的强制前置条件。

TOTP 支持计数防重放与一次性恢复码，登录挑战窗口为五分钟。TOTP 相关秘密依赖服务端主密钥派生与绑定状态；主密钥需要持久化并妥善备份，轮换可能使已有绑定失效。这些控制不包含实名、活体或真实银行卡验证。

A/B/C 为共享演示身份，同一演示 persona 的访客复用同一用户空间。注册账户之间的数据隔离，不意味着共享演示账号的每位访客都独立隔离；同一演示用户的模型与供应商绑定也可能共享。真实授权应使用本人注册账户，并清楚告知数据接收方。

\subsection{接入校验与权限边界}
网关检查 Host、Origin、Sec-Fetch-Site、方法、路径、JSON 类型与请求大小，并设置缓存、内容类型、引用来源和 CSP 等响应策略。公开健康检查及部分认证入口不要求既有会话，业务资源则由会话绑定用户。生产 HTTPS 由 Caddy 终止，Python 接入层提供 HTTP 与应用校验。

\begin{table}[htbp]\centering\small
\caption{主要威胁与控制落点}
\begin{tabularx}{\textwidth}{@{}L{2.5cm}YY@{}}\toprule
威胁 & 现有控制 & 保证范围 \\\midrule
跨用户读取或修改 & 会话绑定用户、对象查询带用户条件 & 注册用户资源；共享 persona 不是访客级隔离 \\
模型越权与注入 & 不可信输入隔离、动作及字段检查、执行网关 & 不代表能识别所有自然语言攻击 \\
过期方案与并发变化 & 修订号、生活版本、TTL、事务内复核 & 契约依业务入口而异 \\
重复扣款与网络重试 & 用户范围唯一键、原方案交易查询、原回执 & 本地合成账本，不外推到供应商清算 \\
凭据泄漏 & 摘要存储、密钥不回显、状态目录隔离 & 用户输入、供应商缓存与日志仍需单独治理 \\\bottomrule
\end{tabularx}
\end{table}

\subsection{可信本机配置与残余风险}
\code{--model-admin} 是可信本机演示开关，不是经过独立管理员角色认证的权限系统。开启后不应让不可信用户访问；公网运行应通过部署环境管理平台密钥。个人模型服务地址的配置校验，也不能证明目标服务值得信任。

反向代理会改写部分请求头，后端看到的 Host 和 Origin 不一定是浏览器原始值，因此应结合实际代理策略评估跨站防护，不能仅列后端头部检查就宣称链路全程保留了原始来源证明。Caddy 可记录请求日志，后端不打印路径不等于整个链路没有日志。真实地址、认证回调、查询参数和凭据缓存都需要适当的访问、保留及清理策略。

\sectionstart
\section{部署拓扑与运行维护}
\subsection{运行环境与服务入口}
后端声明最低 Python 3.10，主体使用标准库；这不表示每个受支持小版本均已完成同等测试。项目 Dockerfile 使用 Python 3.13，并从 Node.js 22 镜像复制 Node 可执行文件。Node 用于官方跑腿及授权组件，也用于前端测试，不是 Python 银行业务服务的替代运行时。

本次文档核验使用 macOS、Python 3.12.14 和 Node.js 24.22.1。它与 Docker 构建声明不同，应分别记录；未提供镜像摘要或容器内版本实测时，不以镜像主版本宣称已完成目标容器验收。

\begin{lstlisting}
./scripts/run.sh
# 或在项目根目录运行
python3 -m backend.server --host 127.0.0.1 --port 8765 \
  --db .local/tropos.sqlite3
\end{lstlisting}
首页为 \code{index.html}，完整工作台为 \code{demo.html}。局域网演示需要显式配置访问方式；设备互访还取决于防火墙与网络隔离。开放演示模式必须视为可修改共享合成状态的环境，不能作为生产访问控制。

\subsection{Linux Docker 与 HTTPS}
当前 Compose 中应用与 Caddy 都使用 host 网络。应用监听宿主机回环 \code{127.0.0.1:8765}，Caddy 在同一宿主网络上接收站点流量并反向代理该地址；配置中没有额外的应用 \code{ports} 发布映射。该拓扑面向 Linux 主机，不能直接推定其他操作系统具有完全相同的 host 网络行为。

应用使用 UID 10001 非 root 用户、只读根文件系统、\code{no-new-privileges} 和受限临时目录。数据库、官方组件及授权状态通过明确目录挂载；组件只读与授权缓存可写需要分别配置权限。配置存在不代表目标服务器的目录属主、域名证书和外网端口已经验收。

\begin{table}[htbp]\centering\small
\caption{关键部署配置及缺失时的行为}
\begin{tabularx}{\textwidth}{@{}L{1.6cm}L{5.1cm}Y@{}}\toprule
配置组 & 配置入口 & 缺失或错误时的影响 \\\midrule
站点与邮件 & \code{TROPOS_SITE_ADDRESS}、\code{TROPOS_PUBLIC_URL}、\code{TROPOS_SMTP_*} & 域名、回调或邮件功能不可据启动成功推定就绪 \\
认证主密钥 & \code{TROPOS_AUTH_MASTER_KEY} 及本地持久化位置 & 自动生成只适于明确管理的环境，丢失或轮换影响 MFA \\
模型服务 & \code{TROPOS_DEEPSEEK_*} 等各服务配置 & 服务不可选或规则降级；真实调用仍依赖供应商权限 \\
美团组件 & \code{TROPOS_MEITUAN_*} & 官方链路不可用，Mock 可独立运行 \\\bottomrule
\end{tabularx}
\end{table}

\subsection{升级与故障恢复}
3.2.x 提供基础 Compose 与版本覆盖文件。更新前备份数据库、配置、认证主密钥与供应商状态，核对挂载目录、UID 和目标镜像；更新后检查健康版本、登录、读取、草案、确认与原回执查询。数据库结构兼容与状态迁移应按实际升级路径验证，不能承诺任何旧版本都能无迁移替换。

\begin{lstlisting}
curl -fsS http://127.0.0.1:8765/api/health
\end{lstlisting}
健康接口返回 \code{version} 为 \code{3.2.9} 只说明服务版本与基础可访问性，不代表 SMTP、真实模型、美团授权和支付均可用。故障时先区分只读请求、已创建草案、已提交订单和结果未知的交易，再按原凭证查询；不以重启或再次提交代替状态核对。

\sectionstart
\section{验证方法与本次核验结果}
\label{sec:validation}
\subsection{可复现基线与计数口径}
本次以 2026 年 11 月 2 日本地工作区复跑自动化测试。核心后端提交为 \code{2a33063}，界面包含尚未提交的水墨风更新，因此结果归属于工作区快照，而非声称某个完全干净的发布提交已经在全部环境验收。源码 SHA-256、运行命令、解释器版本和日志随 LaTeX 源码包的 \code{evidence/} 提供。

\begin{table}[htbp]\centering\small
\caption{本次本机复跑结果}
\begin{tabularx}{\textwidth}{@{}L{2.25cm}L{3.45cm}Y@{}}\toprule
检查层 & 实际计数 & 解释与范围 \\\midrule
Python unittest & 共 470；成功 469；跳过 1；失败与错误 0 & 包括身份、账本、模型替身、HTTP、生活与连接器等测试 \\
Node 测试运行器 & 101 项；通过 101；失败 0；跳过 0 & \code{tests/*.cjs}；并不等于独立的真实浏览器用户测试 \\
规则模式红队 & 攻击 62/62；正常对照 14/14 & 对话攻击 51 项、API 绕过 11 项；当前工作区复跑 \\\bottomrule
\end{tabularx}
\end{table}

Python 跳过项为 \code{test_windows_does_not_apply_unix_chmod_to_credential_directories}，原因是仅验证 Windows 的凭据目录行为，本次 macOS 不适用。跳过不是通过；本次结果也不替代 Windows 实机验证。470 是 unittest 发现并运行的测试总数，已经包含跳过项，不能写成“470 成功再加 1 跳过”。

本次数字按所列命令与日志统计。历史发布材料中的数字需要结合各自日期、环境和计数单位阅读，不能直接替代当前结果。不同层的用例可能重叠，红队场景与 Python 回归也有交叉，不合计为独立覆盖总量。

\begin{lstlisting}
python3 -m unittest discover -s tests -p 'test_*.py' -v
node --test tests/*.cjs
\end{lstlisting}

\subsection{红队样例与结论边界}
锁定样例集包含开发者构造的直接及间接注入、身份冒充、错误完成声明、收款人混淆、金额异常、多步夹带和 API 绕过等场景。攻击通过指不产生攻击者要求的可授权写动作或真实执行，账本保持预期状态且回复不虚称完成；正常对照通过指预期写动作以待授权形式出现，不能把“对照可办理”误写成已经付款。

本次通过 \code{eval.redteam.benchmark(live_model=False)} 在规则模式复跑，将结果单独保存，不覆盖仓库历史报告。历史 \code{reports/redteam.json} 属于 2026 年 10 月 28 日 \code{7ac1c9c} 基线；本次的工作区结果与历史记录分别保留。未调用真实付费模型，也未创建真实美团订单。

62/62 是已知开发样例的回归结果，不是未知攻击拦截率、统计泛化结论或安全认证。模型适配测试大量使用回环 Mock，能检查协议参数和降级行为，不能证明供应商当前账号权限、网络与真实端到端稳定性。

\subsection{与技术主张对应的证据}
\begin{table}[htbp]\centering\small
\caption{关键主张与核验入口}
\begin{tabularx}{\textwidth}{@{}L{3cm}YY@{}}\toprule
技术主张 & 代码与测试入口 & 尚不能据此推出 \\\midrule
确认前不产生普通转账扣款 & \code{service.py:confirm}；转账与 API 测试 & 所有外部交易均可撤销 \\
同一方案不重复记账 & 交易唯一约束、幂等回执及并发测试 & 供应商清算恰好一次 \\
现金路径可复算 & \code{_cash}、生活规划参考路径测试及本文算例 & 日内风险全覆盖或实际到账保证 \\
注册用户数据分区 & \code{test_auth.py} 与对象归属检查 & 共享演示账号访客彼此隔离 \\
模型失败可降级 & \code{test_model.py} 中回环替身与错误响应测试 & 所有供应商真实调用均已验收 \\\bottomrule
\end{tabularx}
\end{table}

真实手机网络、无障碍用户盲测、供应商真实授权支付、目标 Linux 部署、压力与渗透测试需要分别记录环境、输入、结果和失败。功能测试成功不能替代可用性或生产安全验收。

\sectionstart
\section{版本演进与后续工作}
\subsection{3.2 系列的工程收敛}
3.2 将已有转账、账单与订阅能力连接到可解释财务体检，并补充术语解释、长辈友好显示、模型中心与注册账户。后续迭代集中在三个方面：认证交互与输入一致性，官方跑腿组件的授权、部署与错误处理，以及模型自由对话、来源和超时语义。3.2.9 的主要落点是千问普通对话与连接测试的稳定请求策略。

版本号并不增加已经验证的能力范围。水墨视觉更新改善项目表达；财务算法、交易状态与外部服务仍按本文定义运行。发布前应把前端、后端、配置、数据库兼容与验证证据绑定到同一可追溯快照。

\subsection{优先完成的工程改进}
\begin{enumerate}
\item \textbf{统一动作确认契约。} 将普通转账、多步骤任务、订阅、卡片与理财变更逐一映射到相同的核验清单；尤其补齐执行时的锁定期、卡状态等业务重验，明确各入口的限额、版本和时限，不依赖规划时曾经检查过的条件。
\item \textbf{细化任务终态与恢复。} 将整单执行尝试结束与全部成功区分，面向用户展示部分完成、失败及暂停项；重新办理时引用原回执，避免补做任务导致重复。
\item \textbf{收紧对外数据控制。} 对当前消息、历史、目标名称和检索词提供清楚的发送提示与可测试的字段治理；对真实美团地址、供应商响应和缓存建立保留与访问策略。
\item \textbf{分开演示与真实授权空间。} 共享 persona 用于合成演示，真实外部账号使用独立注册身份。公网配置平台密钥应采用明确管理员身份与最小权限，而不是仅依靠本机开关。
\item \textbf{完善目标环境验证。} 固化 Linux 镜像与官方组件版本，补 Windows、真实手机及外部服务验收，建立备份恢复和关键状态观测。SQLite 单机适用于当前演示；多实例演进需重新设计数据库锁、共享任务与幂等边界。
\end{enumerate}

\subsection{面向更完整金融生活的方向}
项目可以继续探索家庭协同、跨账户资金视图、小微经营中的应收应付与回款压力，以及更长时期的目标衔接。前提是获得相应数据与授权，明确金融业务资质、供应商契约和可验证的决策边界。不能把当前 30 天有限模型外推为已经实现的完整生命周期优化。

TROPOS 的长期方向仍是同一件事：理解人的处境，让安排跟随变化，同时把决定权与可追溯的结果留给用户。技术演进应以这一要求衡量，而不是以增加模型数量或场景数量作为完成标准。

\appendix
\sectionstart
\section{接口与工程核对索引}
\subsection{主要接口契约}
以下列出代表入口，完整字段和错误行为以 \code{backend/server.py} 及对应服务测试为准。公开入口与登录业务入口分开；需要变更状态的请求不因路径存在就绕过用户确认。

\small\setlength{\tabcolsep}{4pt}
\begin{longtable}{@{}L{1.5cm}L{6.9cm}L{5.0cm}@{}}
\caption{代表性 HTTP 接口}\\
\toprule 方法 & 路径 & 用途 \\\midrule\endfirsthead
\multicolumn{3}{l}{\small 表\thetable\quad 代表性 HTTP 接口（续）}\\
\toprule 方法 & 路径 & 用途 \\\midrule\endhead
\bottomrule\endfoot
GET & \code{/api/health} & 公开健康与版本 \\
POST & \code{/api/auth/register}\newline\code{/api/auth/login} & 注册和登录 \\
POST & \code{/api/auth/email/verify}\newline\code{/api/auth/password/reset} & 单次挑战验证与密码重置 \\
GET & \code{/api/state}\newline\code{/api/insights} & 当前快照与财务体检 \\
GET & \code{/api/financial-glossary} & 只读金融术语 \\
POST & \code{/api/agent} & 对话理解与业务分派 \\
GET & \code{/api/conversations/{id}} & 恢复对话与资源状态 \\
POST & \code{/api/transfers/propose} & 生成普通转账草案 \\
POST & \code{/api/proposals/{id}/confirm} & 确认有效草案 \\
POST & \code{/api/missions}\newline\code{/api/missions/{id}/confirm} & 多步骤规划与逐项批准 \\
POST & \code{/api/life/analyze}\newline\code{/api/life/commit}\newline\code{/api/life/goal} & 比较、登记义务、另建目标草案 \\
GET & \code{/api/model-center} & 模型配置状态与能力投影 \\
PUT & \code{/api/model-center/selection} & 当前用户选择模型 \\
POST & \code{/api/shopping/paotui/preview} & 官方帮买预览 \\
POST & \code{/api/shopping/paotui/previews/{id}/confirm} & 逐项确认并提交官方订单 \\
\end{longtable}
\normalsize

\clearpage
\subsection{工程模块与核查函数}
\small
\begin{longtable}{@{}L{4.8cm}L{4.6cm}L{4.0cm}@{}}
\caption{核心实现索引}\\
\toprule 文件 & 代表函数或对象 & 重点核对内容 \\\midrule\endfirsthead
\multicolumn{3}{l}{\small 表\thetable\quad 核心实现索引（续）}\\
\toprule 文件 & 代表函数或对象 & 重点核对内容 \\\midrule\endhead
\bottomrule\endfoot
\code{backend/service.py} & \code{_cash}, \code{_checks}, \code{confirm}, \code{_post} & 现金路径、复核、记账与幂等 \\
\code{backend/conversation.py} & \code{turn}, \code{_draft_transfer}, \code{_refresh} & 对话版本、草案更正和恢复 \\
\code{backend/agent_ops.py} & \code{plan}, \code{confirm}, \code{_execute} & 多步骤部分执行与回执 \\
\code{backend/life_planner.py} & \code{analyze}, \code{_choices}, \code{_rank} & 四情景、离散选择、束宽与排序 \\
\code{backend/life_service.py} & \code{commit_plan}, \code{propose_goal} & 登记义务与目标草案分离 \\
\code{backend/insights.py} & \code{build_insights} & 可见数据统计与解释规则 \\
\code{backend/model.py} & 模型适配、建议验证与本地策略 & 预算、数据发送与失败语义 \\
\code{backend/model_center.py} & 用户选择、配置与能力状态 & 平台配置及个人密钥边界 \\
\code{backend/meituan_paotui.py} & 组件调用与用户目录管理 & 参数数组、缓存与供应商错误 \\
\code{backend/shopping.py} & Mock 与官方订单业务 & 订单未知状态和账本分离 \\
\code{backend/server.py} & HTTP 路由与接入校验 & 会话归属、方法、头部与参数 \\
\code{deploy/docker-compose.yml} & 应用与 Caddy host 网络 & 回环上游与目录挂载 \\
\end{longtable}
\normalsize

\subsection{交付与复现}
主文件为 \code{main.tex}，分节正文在 \code{sections/}，项目 logo 与长卷原件在 \code{assets/}。使用 kit 的中文 paper 日常版，以 XeLaTeX 编译；环境与依赖说明见源码包 README。
\begin{lstlisting}
latexmk -xelatex -interaction=nonstopmode \
  -halt-on-error main.tex
\end{lstlisting}
\code{evidence/} 保存测试日志、规则红队报告、代码指纹与核验说明，用于追溯本次结果。源码包不包含用户数据库、密钥或供应商认证材料。

\end{document}
